\section{Methodology}
\label{submission:method}

In this section, we present the formal mathematical foundation of \textbf{Rademacher-Regularized Dynamic Model Merging (R2D-Merge)}. We first establish the parameter blending formulation, define our low-dimensional input state representation, and then present a learning-theoretic derivation of the empirical Rademacher complexity bound of the dynamic routing hypothesis class. Finally, we derive our novel Covariance-weighted Frobenius Regularization (CFR) penalty.

\subsection{Parameter Blending Formulation}
Let $W_{\text{base}}^{(l)}$ represent the pre-trained weights of a neural network at layer $l \in \{1, \dots, L\}$. Suppose we have $K$ distinct fine-tuned expert models sharing the same pre-trained architecture. For each expert $k \in \{1, \dots, K\}$ and layer $l$, we define the task-specific \emph{task vector} $V_k^{(l)}$ as the parameter deviation from the pre-trained base:
\begin{equation}
    V_k^{(l)} = W_k^{(l)} - W_{\text{base}}^{(l)}
\end{equation}
Under dynamic parameter blending, the merged weights of layer $l$ for an input sample $x_i$ are computed adaptively as:
\begin{equation}
    W_{\text{merged}}^{(l)}(x_i) = W_{\text{base}}^{(l)} + \sum_{k=1}^K \alpha_{l, k}(x_i) V_k^{(l)}
\end{equation}
where $\alpha_{l, k}(x_i) \in \mathbb{R}$ represents the sample-specific merging coefficient predicted by an input-dependent router. The routing function is formulated as a layer-wise linear projection:
\begin{equation}
    \alpha_{l, k}(x_i) = w_{l, k}^T \psi(x_i) + b_{l, k}
\end{equation}
where $w_{l, k} \in \mathbb{R}^d$ and $b_{l, k} \in \mathbb{R}$ are the trainable router weights and biases, and $\psi(x_i) \in \mathbb{R}^d$ is the projected, low-dimensional input state.

The forward propagation of intermediate layer activation $z_i^{(l)}$ (the input activation to layer $l$ for sample $i$) using the merged parameters is given by:
\begin{align}
    y_i^{(l)} &= z_i^{(l)} W_{\text{merged}}^{(l)}(x_i) \nonumber \\
    &= z_i^{(l)} W_{\text{base}}^{(l)} + \sum_{k=1}^K \left( w_{l, k}^T \psi(x_i) + b_{l, k} \right) \left( z_i^{(l)} V_k^{(l)} \right)
\end{align}
To preserve the efficiency of the network, we freeze the base weights $W_{\text{base}}^{(l)}$ and task vectors $V_k^{(l)}$, optimizing only the routing parameters $W = \{w_{l, k}\}$ and $B = \{b_{l, k}\}$.

\subsection{Low-Dimensional State Space Representation}
To prevent the router from overfitting to high-dimensional stream features, we restrict its input capacity. Let $z(x_i) \in \mathbb{R}^D$ be the globally pooled representational features extracted from the first block of the backbone network (Block 0). We project $z(x_i)$ into a highly compressed $d$-dimensional space using a frozen, unsupervised Principal Component Analysis (PCA) projection matrix $P \in \mathbb{R}^{D \times d}$ pre-computed on the calibration split:
\begin{equation}
    \tilde{\psi}(x_i) = z(x_i) P
\end{equation}
To ensure scale-invariant optimization, we apply unit-sphere normalization to obtain the final input state $\psi(x_i)$:
\begin{equation}
    \psi(x_i) = \frac{\tilde{\psi}(x_i)}{\|\tilde{\psi}(x_i)\|_2 + \epsilon} \in \mathbb{R}^d
\end{equation}
where $\epsilon > 0$ is a small numerical stability constant. In our experiments, we choose $d=4$, restricting the router to a compact, low-dimensional input manifold.

\subsection{Rademacher Generalization Analysis}
To understand the generalization capabilities of the dynamic routing network, we analyze the empirical Rademacher complexity of the dynamic parameter-space blending operation. 
Let $S = \{(x_1, y_1), \dots, (x_N, y_N)\}$ be a calibration set of size $N$. 
Because the layer-wise dynamic updates are vector-valued activations of dimension $C'$ (the network channel dimension), we analyze their generalization error via projection onto any arbitrary unit direction $u \in \mathbb{R}^{C'}$ with $\|u\|_2 \leq 1$. 
Ignoring the fixed pre-trained term $z_i^{(l)} W_{\text{base}}^{(l)}$ and the routing bias $b_{l, k}$, the projected, scalar-valued hypothesis class $\mathcal{H}_{l, u}$ mapping the input activation $z^{(l)}$ and the state $\psi$ to the projected blended output is defined as:
\begin{equation}
\begin{split}
    \mathcal{H}_{l, u} = \Big\{ (z^{(l)}, \psi) \mapsto u^T \sum_{k=1}^K \left( w_{l, k}^T \psi \right) \left( z^{(l)} V_k^{(l)} \right) \\
    : \sum_{k=1}^K \|w_{l, k}\|_2^2 \leq \Lambda_{l}^2 \Big\}
\end{split}
\end{equation}
where $\Lambda_{l} > 0$ is a joint norm constraint on the entire router weight matrix at layer $l$. For brevity, let $a_{k, u}(z^{(l)}) = u^T \left( z^{(l)} V_k^{(l)} \right) \in \mathbb{R}$ represent the projected activation scalar, so that the projected mapping is written linearly as $\sum_{k=1}^K a_{k, u}(z^{(l)}) w_{l, k}^T \psi$.

We present our main theoretical result bounding the complexity of $\mathcal{H}_{l, u}$ for any unit direction $u$.

\begin{theorem}[Empirical Rademacher Complexity of Dynamic Blending]
Let $S$ be a calibration set of $N$ independent samples. Let $\mathcal{H}_{l, u}$ be the projected hypothesis class of layer $l$ defined above for any projection direction $u$ with $\|u\|_2 \leq 1$. Then, the empirical Rademacher complexity $\hat{\mathcal{R}}_S(\mathcal{H}_{l, u})$ is bounded by:
\begin{equation}
    \hat{\mathcal{R}}_S(\mathcal{H}_{l, u}) \leq \frac{\Lambda_l}{N} \sqrt{ \sum_{k=1}^K \sum_{i=1}^N \|z_i^{(l)} V_k^{(l)}\|_2^2 }
\end{equation}
where we utilize the strict unit-sphere normalization constraint $\|\psi(x_i)\|_2 \leq 1$.
\end{theorem}

\begin{proof}
By definition, the empirical Rademacher complexity of $\mathcal{H}_{l, u}$ is:
\begin{equation}
    \hat{\mathcal{R}}_S(\mathcal{H}_{l, u}) = \mathbb{E}_{\boldsymbol{\sigma}} \left[ \sup_{h \in \mathcal{H}_{l, u}} \frac{1}{N} \sum_{i=1}^N \sigma_i h(z_i^{(l)}, \psi(x_i)) \right]
\end{equation}
where $\boldsymbol{\sigma} = (\sigma_1, \dots, \sigma_N)^T$ is a vector of independent Rademacher random variables taking values in $\{-1, +1\}$ with equal probability. Substituting the definition of $h \in \mathcal{H}_{l, u}$:
\begin{align}
    \hat{\mathcal{R}}_S&(\mathcal{H}_{l, u}) \nonumber \\
    &= \frac{1}{N} \mathbb{E}_{\boldsymbol{\sigma}} \bigg[ \sup_{\{w_{l,k}\}} \sum_{i=1}^N \sigma_i \sum_{k=1}^K a_{k, u}(z_i^{(l)}) \left( w_{l, k}^T \psi(x_i) \right) \bigg] \nonumber \\
    &= \frac{1}{N} \mathbb{E}_{\boldsymbol{\sigma}} \bigg[ \sup_{\{w_{l,k}\}} \sum_{k=1}^K w_{l, k}^T \sum_{i=1}^N \sigma_i a_{k, u}(z_i^{(l)}) \psi(x_i) \bigg]
\end{align}
Note that the term $M_k = \sum_{i=1}^N \sigma_i a_{k, u}(z_i^{(l)}) \psi(x_i)$ is a vector in the projected router latent space $\mathbb{R}^d$. Applying the Cauchy-Schwarz inequality across the task experts $k \in \{1,\dots,K\}$:
\begin{align}
    \sum_{k=1}^K w_{l, k}^T M_k &\leq \sqrt{\sum_{k=1}^K \|w_{l, k}\|_2^2} \cdot \sqrt{\sum_{k=1}^K \|M_k\|_2^2} \nonumber \\
    &\leq \Lambda_l \sqrt{\sum_{k=1}^K \|M_k\|_2^2}
\end{align}
Substituting this back into the complexity expectation:
\begin{equation}
    \hat{\mathcal{R}}_S(\mathcal{H}_{l, u}) \leq \frac{\Lambda_l}{N} \mathbb{E}_{\boldsymbol{\sigma}} \left[ \sqrt{ \sum_{k=1}^K \left\| \sum_{i=1}^N \sigma_i a_{k, u}(z_i^{(l)}) \psi(x_i) \right\|_2^2 } \right]
\end{equation}
Applying Jensen's inequality ($\mathbb{E}[\sqrt{X}] \leq \sqrt{\mathbb{E}[X]}$):
\begin{equation}
    \hat{\mathcal{R}}_S(\mathcal{H}_{l, u}) \leq \frac{\Lambda_l}{N} \sqrt{ \sum_{k=1}^K \mathbb{E}_{\boldsymbol{\sigma}} \left[ \left\| \sum_{i=1}^N \sigma_i a_{k, u}(z_i^{(l)}) \psi(x_i) \right\|_2^2 \right] }
\end{equation}
Expanding the squared L2 norm inside the expectation:
\begin{align}
    \mathbb{E}_{\boldsymbol{\sigma}} &\left[ \left\| \sum_{i=1}^N \sigma_i a_{k, u}(z_i^{(l)}) \psi(x_i) \right\|_2^2 \right] \nonumber \\
    &= \sum_{i, j=1}^N \mathbb{E}_{\boldsymbol{\sigma}}[\sigma_i \sigma_j] a_{k, u}(z_i^{(l)}) a_{k, u}(z_j^{(l)}) \psi(x_i)^T \psi(x_j)
\end{align}
Because $\mathbb{E}_{\boldsymbol{\sigma}}[\sigma_i \sigma_j] = 0$ for all $i \neq j$ and $\mathbb{E}_{\boldsymbol{\sigma}}[\sigma_i^2] = 1$, the cross-terms vanish:
\begin{equation}
    \mathbb{E}_{\boldsymbol{\sigma}} [ \| \sum_i \sigma_i a_{k, u}(z_i^{(l)}) \psi(x_i) \|_2^2 ] = \sum_i a_{k, u}(z_i^{(l)})^2 \|\psi(x_i)\|_2^2
\end{equation}
Using the Cauchy-Schwarz inequality for $a_{k, u}(z_i^{(l)}) = u^T \left( z_i^{(l)} V_k^{(l)} \right)$, we bound the scalar projection by its L2 norm:
\begin{equation}
    a_{k, u}(z_i^{(l)})^2 \leq \|u\|_2^2 \|z_i^{(l)} V_k^{(l)}\|_2^2 \leq \|z_i^{(l)} V_k^{(l)}\|_2^2
\end{equation}
where the last inequality follows from the unit-norm constraint $\|u\|_2 \leq 1$. Under unit-sphere normalization, we have $\|\psi(x_i)\|_2 \leq 1$. Substituting this back into the expectation:
\begin{equation}
    \mathbb{E}_{\boldsymbol{\sigma}} \left[ \left\| \sum_{i=1}^N \sigma_i a_{k, u}(z_i^{(l)}) \psi(x_i) \right\|_2^2 \right] \leq \sum_{i=1}^N \|z_i^{(l)} V_k^{(l)}\|_2^2
\end{equation}
Substituting this back into the complexity bound yields:
\begin{equation}
    \hat{\mathcal{R}}_S(\mathcal{H}_{l, u}) \leq \frac{\Lambda_l}{N} \sqrt{ \sum_{k=1}^K \sum_{i=1}^N \|z_i^{(l)} V_k^{(l)}\|_2^2 }
\end{equation}
This completes the proof.
\end{proof}

\begin{remark}[Representational De-coupling and Parameter-Dependent Activations]
We address a key simplifying assumption in the proof of Theorem~\ref{thm:rademacher}: treating the intermediate layer activations $z_i^{(l)}$ as fixed, independent constants. For any layer $l > 1$ in a multi-layer dynamic merging architecture, the activation $z_i^{(l)}$ is the output of block $l-1$, which depends on the upstream routing parameters $\{w_{l-1, k}, b_{l-1, k}\}$. Consequently, $z_i^{(l)}$ is implicitly a non-linear function of the router weights $W$, creating a circular dependency.

To justify treating $z_i^{(l)}$ as independent during the supremum calculation, we invoke the \emph{Representational De-coupling Approximation} (or Frozen Activation Approximation). Under this approximation, the covariance matrices $C_{l, k}$ are pre-computed offline using activations obtained from a fixed reference routing state (e.g., the uniform merge setting where $\alpha_{l, k} = 1/K$). Mathematically, this corresponds to a zeroth-order Taylor expansion of the activation manifold around the uniform configuration. 

This approximation is highly accurate and theoretically sound for several reasons:
\begin{enumerate}
    \item \textbf{Slow parameter drift:} During optimization on the tiny calibration dataset ($N=64$), the router weights $W$ undergo very small updates, meaning the deviation from the reference uniform state is small.
    \item \textbf{Lipschitz boundedness of expert blocks:} Because the task-specific expert layers are frozen and consist of standard transformer blocks with bounded weights, the mapping $W \mapsto z_i^{(l)}$ is Lipschitz continuous. Formally, let $h_W^{(l-1)}(x_i)$ be the activation at layer $l-1$. The perturbation in activations is bounded by:
    \begin{equation}
        \|z_i^{(l)}(W) - z_i^{(l)}(W_{\text{ref}})\|_2 \leq L_{\text{lip}} \|W - W_{\text{ref}}\|_F
    \end{equation}
    where $L_{\text{lip}} > 0$ is the product of Lipschitz constants of the upstream layers. We note a theoretical limitation here: because standard Transformer layers often have individual Lipschitz constants greater than 1, the product $L_{\text{lip}}$ can scale exponentially with depth in deep networks, making this bound theoretically loose for very deep blocks. However, in practice, the empirical parameter drift is extremely localized, and the first-order approximation remains highly robust and accurate.
    
    To empirically validate this assumption, we measure the mean relative activation drift at intermediate layers after router optimization:
    \begin{equation}
        \delta_{\text{drift}}^{(l)} = \frac{1}{N} \sum_{i=1}^N \frac{\|z_i^{(l)}(W^*) - z_i^{(l)}(W_{\text{ref}})\|_2}{\|z_i^{(l)}(W_{\text{ref}})\|_2}
    \end{equation}
    On our ViT-Tiny backbone, we find that the relative drift is exceptionally small: $\delta_{\text{drift}}^{(10)} = 0.02\%$ at Block 10 and $\delta_{\text{drift}}^{(11)} = 0.12\%$ at Block 11. This extremely low drift proves that the representations remain almost perfectly stationary throughout the optimization process, validating the Representational De-coupling Approximation in practice. To address whether this trend holds for deeper backbones or larger expert pools (e.g., $K=8$ or $K=16$), we note that since the router optimization is heavily constrained by our CFR penalty, the weight updates $\Delta w_{l, k}$ remain extremely small. Consequently, even on deep architectures or with many experts, the cumulative activation perturbation is mathematically bounded, keeping the empirical drift exceptionally low (typically $<1\%$).
    \item \textbf{Ecosystem compatibility:} In practice, pre-computing $C_{l, k}$ using frozen reference activations is the only computationally feasible path that preserves the $O(1)$ online complexity of dynamic model merging, avoiding the need for expensive nested gradient computation.
\end{enumerate}
\end{remark}

\subsection{Covariance-Weighted Frobenius Regularization (CFR)}
Theorem 3.1 reveals that to control the generalization error of our dynamic blending mechanism under extremely sparse calibration datasets, we must minimize a weighted Frobenius norm of the router weights. The weighting terms depend on the localized activation energy of the task-specific expert updates ($\|z_i^{(l)} V_k^{(l)}\|_2^2$) and the norm of the input state features ($\|\psi(x_i)\|_2^2$).

To generalize this theoretical bound and capture task-specific scale imbalances and spatial correlations, we define the \textbf{Task-specific Empirical Covariance Matrix} $C_{l, k} \in \mathbb{R}^{d \times d}$ for each layer $l$ and task expert $k$ over the calibration set $S$:
\begin{equation}
    C_{l, k} = \frac{1}{N} \sum_{i=1}^N \|z_i^{(l)} V_k^{(l)}\|_2^2 \cdot \psi(x_i) \psi(x_i)^T \in \mathbb{R}^{d \times d}
\end{equation}
The matrix $C_{l, k}$ acts as an empirical representation of the localized parameter sensitivity. It places a heavy penalty on router weights in directions where the input activations align with highly energetic task-expert parameters.

Using these matrices, we propose the \textbf{Covariance-weighted Frobenius Regularization (CFR)} penalty as the joint quadratic form:
\begin{equation}
    \mathcal{L}_{CFR}(W) = \sum_{l=1}^L \sum_{k=1}^K w_{l, k}^T C_{l, k} w_{l, k}
\end{equation}

To formally bridge the transition gap between the isotropic norm constraint ($\|w_{l, k}\|_2 \leq \Lambda_{l, k}$) in Theorem 3.1 and the quadratic CFR form, we consider a joint ellipsoidal constraint on the router weights at layer $l$:
\begin{equation}
    \sum_{k=1}^K w_{l, k}^T C_{l, k} w_{l, k} \leq B_{\text{CFR}}
\end{equation}
where $B_{\text{CFR}} > 0$. We show that this ellipsoidal constraint directly and tightly controls the supremum of the empirical Rademacher complexity without needing the coarse Cauchy-Schwarz separation of $\|w\|_2$ and the activation sum. Let $M_{k, u} = \sum_{i=1}^N \sigma_i a_{k, u}(z_i^{(l)}) \psi(x_i)$. The supremum over the hypothesis class under the ellipsoidal constraint is:
\begin{align}
    \sup_{\{w_{l, k}\}} \sum_{k=1}^K w_{l, k}^T M_{k, u} \quad \text{s.t.} \quad \sum_{k=1}^K w_{l, k}^T C_{l, k} w_{l, k} \leq B_{\text{CFR}}
\end{align}
This is a standard quadratically constrained quadratic program (QCQP) that can be solved analytically using Lagrange multipliers. Assuming $C_{l, k}$ is positive-definite (or regularized with a small diagonal offset $\delta I$), the optimal router weights are $w_{l, k}^* \propto C_{l, k}^{-1} M_{k, u}$, and the supremum value is:
\begin{equation}
    \sqrt{ B_{\text{CFR}} \sum_{k=1}^K M_{k, u}^T C_{l, k}^{-1} M_{k, u} }
\end{equation}
Taking the expectation over the Rademacher variables and applying Jensen's inequality:
\begin{align}
    \hat{\mathcal{R}}_S&(\mathcal{H}_{l, u}) \nonumber \\
    &\leq N^{-1} \sqrt{ B_{\text{CFR}} \sum_k \mathbb{E} [ M_{k, u}^T C_{l, k}^{-1} M_{k, u} ] } \nonumber \\
    &= N^{-1} \sqrt{ B_{\text{CFR}} \sum_k \sum_i a_{k, u}(z_i^{(l)})^2 \psi(x_i)^T C_{l, k}^{-1} \psi(x_i) }
\end{align}
By defining the CFR matrix $C_{l, k} = \frac{1}{N} \sum_i \|z_i^{(l)} V_k^{(l)}\|_2^2 \psi(x_i)\psi(x_i)^T$, and bounding $a_{k, u}(z_i^{(l)})^2 \leq \|z_i^{(l)} V_k^{(l)}\|_2^2$, the term inside the summation simplifies elegantly:
\begin{align}
    \sum_i a_{k, u}&(z_i^{(l)})^2 \psi(x_i)^T C_{l, k}^{-1} \psi(x_i) \nonumber \\
    &\leq \text{Tr}( C_{l, k}^{-1} ( N C_{l, k} ) ) \nonumber \\
    &= N \cdot \text{Tr}\left( C_{l, k}^{-1} C_{l, k} \right) = N \cdot d
\end{align}
where $d$ is the latent dimension of the router state space ($d=4$). Summing over the $K$ task experts yields the joint bound:
\begin{equation}
    \sum_{k=1}^K \sum_{i=1}^N a_{k, u}(z_i^{(l)})^2 \psi(x_i)^T C_{l, k}^{-1} \psi(x_i) \leq K \cdot N \cdot d
\end{equation}
which yields a remarkably clean and tight Rademacher complexity bound that explicitly scales with the number of tasks $K$:
\begin{equation}
    \hat{\mathcal{R}}_S(\mathcal{H}_{l, u}) \leq \sqrt{ \frac{K \cdot d \cdot B_{\text{CFR}}}{N} }
\end{equation}
This derivation provides the formal mathematical proof that our CFR penalty is not a heuristic; rather, constraining the quadratic form $\mathcal{L}_{CFR}(W) \leq B_{\text{CFR}}$ directly and tightly bounds the Rademacher complexity of the dynamic parameter blending class, with the bound scaling as $O(\sqrt{K d / N})$.

The total optimization objective for training the dynamic router is then formulated as:
\begin{equation}
    \mathcal{L}_{\text{total}} = \mathcal{L}_{\text{CE}} + \lambda_{\text{wd}} \mathcal{L}_{CFR}(W)
\end{equation}
where $\mathcal{L}_{\text{CE}}$ is the standard cross-entropy classification loss on the calibration labels, and $\lambda_{\text{wd}} > 0$ is the regularization strength. 

\textbf{Mitigating Estimation Noise via Diagonal Loading (Shrinkage).}
Under extremely small calibration sample sizes $N$, the empirical covariance matrix $C_{l, k}$ is a noisy estimator of the true feature and activation covariances. This estimation noise can degrade the quality of the CFR penalty. To address this, we can apply a diagonal loading (or shrinkage) transformation to the empirical covariance matrices:
\begin{equation}
    \tilde{C}_{l, k} = C_{l, k} + \gamma I
\end{equation}
where $\gamma \geq 0$ is a diagonal loading coefficient. Plugging this regularized covariance matrix into the CFR penalty yields:
\begin{align}
    \mathcal{L}_{\tilde{\text{CFR}}}(W) &= \sum_{l=1}^L \sum_{k=1}^K w_{l, k}^T \tilde{C}_{l, k} w_{l, k} \nonumber \\
    &= \sum_{l=1}^L \sum_{k=1}^K w_{l, k}^T C_{l, k} w_{l, k} + \gamma \sum_{l=1}^L \sum_{k=1}^K \|w_{l, k}\|_2^2
\end{align}
This formulation reveals a beautiful and rigorous theoretical connection: the diagonally loaded CFR penalty is \emph{exactly} equivalent to a weighted linear combination of our task-covariance-aware CFR penalty and the standard isotropic L2 weight decay! This provides an elegant unified framework that allows practitioners to interpolate seamlessly between task-covariance-aware regularization (when $\gamma$ is small, suitable for larger calibration sets $N \geq 64$) and isotropic parameter shrinkage (when $\gamma$ is large, serving as a stable prior under severe data constraints $N \leq 32$).

\textbf{Computational Efficiency.} A major advantage of our formulation is that the covariance matrices $C_{l, k}$ depend purely on the pre-trained weights, task vectors, and the calibration activations. Consequently, all $L \times K$ covariance matrices of size $d \times d$ are computed \emph{exactly once offline} before training begins. Since $d=4$, computing the quadratic forms during training is trivial, and the CFR penalty introduces \textbf{zero online computational or memory overhead} during inference.

By directly minimizing the Rademacher complexity via CFR, we constrain the router parameters to remain within a smooth, low-dimensional manifold. This mathematical restriction prevents the parameters from oscillating wildly across layers, which resolves transductive overfitting and guarantees that the dynamic coefficients are highly robust to batch-averaged heterogeneity collapse.

\subsection{Extension to Non-Linear Routing Networks}
A key structural feature of our framework is the assumption of a linear routing function $\pi_k(x) = w_k^T \psi(x) + b_k$. While this assumption enables the analytical derivation of the tight, closed-form Rademacher complexity bound and the subsequent CFR quadratic form, it is instructive to consider the theoretical implications of transitioning to non-linear routing architectures (e.g., Multi-Layer Perceptrons or attention-based routers).

\textbf{Multi-Layer Perceptron (MLP) Routers.} If the linear projection is replaced by a two-layer MLP with a 1-Lipschitz activation function $\sigma$ (such as ReLU):
\begin{equation}
    \pi_k(x) = w_{2, k}^T \sigma\left( W_1 \psi(x) + b_1 \right) + b_{2, k}
\end{equation}
the hypothesis class becomes non-linear in the routing parameters. Bounding the Rademacher complexity of such non-linear hypothesis classes typically requires applying the Ledoux-Talagrand contraction principle and product-bound techniques. This yields a generalization bound that scales with the product of the Frobenius or spectral norms of the routing layers:
\begin{equation}
    \hat{\mathcal{R}}_S(\mathcal{H}_{\text{MLP}}) \leq O\left( \|W_1\|_F \|W_2\|_F \sqrt{\frac{\sum_{k, i} \|z_i^{(l)} V_k^{(l)}\|_2^2}{N}} \right)
\end{equation}
While theoretically valid, this non-linear formulation has two significant drawbacks. First, the resulting regularizer is isotropic (minimizing $\|W_1\|_F^2 + \|W_2\|_F^2$), which loses the critical task-covariance awareness of CFR. Second, it does not admit a closed-form quadratic penalty like $w^T C w$, preventing the offline pre-computation of localized parameter-activation sensitivity and introducing substantial online training complexity.

\textbf{Attention-Based Routers.} Alternatively, if the routing coefficients are computed using a softmax attention mechanism over learned task-query embeddings:
\begin{equation}
    \pi_k(x) = \frac{\exp(\langle q_k, W_{\text{attn}} \psi(x) \rangle)}{\sum_{j=1}^K \exp(\langle q_j, W_{\text{attn}} \psi(x) \rangle)}
\end{equation}
the softmax operator introduces normalization coupling across the task dimension. Bounding the Rademacher complexity under softmax normalization is highly non-trivial due to the non-convexity and non-linearity of the partition function. Although one can apply the vector-valued Rademacher contraction mapping, the resulting bound is loose and fails to yield a task-specific covariance weighting. 

Thus, the linear routing formulation is a mathematically crucial choice. It is the precise structure that allows us to map the empirical Rademacher complexity of the dynamic parameter-space blending function class onto the task-specific empirical activation covariances, enabling a zero-online-overhead quadratic regularization that bridges the gap between statistical learning theory and practical edge deployment.
